\section{实验五：基于颜色识别的视觉抓取}

\subsection{实验目的}

本实验在相机标定代码基础上，通过颜色识别获得物体在图像中的位置，再将像素坐标转换为机械臂基坐标，并生成抓取与放置轨迹的仿真方案。相比前面的参考点回放，本实验由视觉结果计算候选抓取点，体现了从固定动作到视觉引导规划的程序结构转变；本组未执行真实抓取。

\subsection{颜色分割与目标定位}

颜色识别代码主要位于 \path{daran/color_pick_place/detect.py}。程序定义了红、绿、蓝、黄四类颜色的 HSV 阈值范围，对图像进行颜色分割后，通过轮廓提取获得候选色块。检测流程如下：

\begin{enumerate}
  \item 读取相机图像并根据标定参数进行去畸变。
  \item 将 BGR 图像转换到 HSV 颜色空间。
  \item 根据目标颜色阈值生成二值 mask。
  \item 使用开运算和闭运算去除噪声、填补孔洞。
  \item 通过 \texttt{cv2.findContours} 提取轮廓。
  \item 根据面积阈值和长宽比过滤非方块目标。
  \item 使用 \texttt{cv2.minAreaRect} 得到目标中心、边界框和旋转角。
\end{enumerate}

对于蓝色和绿色，程序还叠加了 BGR 通道约束。例如蓝色目标要求蓝色通道显著高于红色和绿色通道，以增强在复杂光照下的鲁棒性。

各颜色的 HSV 阈值如表 \ref{tab:color-thresholds} 所示。红色在 HSV 中跨越 0 度边界，因此需要使用两个区间合并；蓝色和绿色除 HSV 阈值外还加入了原始 BGR 通道约束，用于过滤偏灰或反光区域。

\begin{table}[H]
  \centering
  \caption{颜色识别 HSV 阈值设置}
  \label{tab:color-thresholds}
  \small
  \begin{tabularx}{0.88\textwidth}{L{1.6cm}Y}
    \toprule
    颜色 & HSV 范围 \\
    \midrule
    red & $[0,100,100]\sim[10,255,255]$，$[160,100,100]\sim[179,255,255]$ \\
    green & $[35,50,50]\sim[85,255,255]$ \\
    blue & $[78,70,35]\sim[135,255,255]$ \\
    yellow & $[20,100,100]\sim[30,255,255]$ \\
    \bottomrule
  \end{tabularx}
\end{table}

选择 HSV 颜色空间的原因是其将色相、饱和度和亮度近似分离。相比直接在 RGB/BGR 空间设阈值，HSV 中的 H 通道更接近人对“颜色类别”的描述，S 和 V 通道则可以用于过滤灰暗区域或高亮反光区域。不过 HSV 方法仍然依赖光照条件，尤其在阴影、曝光变化和物体表面反光时，同一颜色物体的像素分布可能跨出预设阈值范围。

形态学开闭运算用于改善二值 mask 的连通性。开运算可以去除零散噪声点，闭运算可以填补目标区域内部的小孔洞。经过形态学处理后，轮廓提取结果更稳定，\texttt{minAreaRect} 得到的中心点和角度也更适合传递给后续抓取规划模块。

颜色阈值分割和轮廓提取的核心代码如代码 \ref{lst:exp5-color-code} 所示。该代码保留了原实验中“HSV 阈值 + 形态学处理 + 最小外接矩形”的基本流程。

\lstinputlisting[
  language=Python,
  caption={颜色阈值分割与色块轮廓提取},
  label={lst:exp5-color-code}
]{code/exp5_color_visual_grasp.py}

\begin{figure}[H]
  \centering
  \includegraphics[width=0.72\textwidth]{task5_color_raw_image.png}
  \caption{颜色识别实验采集图像示例}
  \label{fig:task5-color-raw}
\end{figure}

\subsection{像素坐标到机械臂坐标转换}

坐标转换由 \path{daran/color_pick_place/convert.py} 实现。检测结果可以是像素点 $(u,v)$，也可以包含角度信息的字典：
\[
\{ \text{center}:(u,v),\ \text{angle}:\theta \}.
\]
转换步骤如下：

\begin{enumerate}
  \item 读取 \path{camera_calibration_result.json} 中的相机内参和畸变参数。
  \item 读取 \path{extrinsics.json} 中相机到机械臂基坐标的齐次变换 $T_{\text{robot}}^{\text{camera}}$。
  \item 使用 \texttt{cv2.undistortPoints} 将像素点转换为相机归一化射线。
  \item 将射线变换到机械臂坐标系，并与桌面平面 $z=\SI{80}{mm}$ 求交。
  \item 输出机械臂坐标 $(x,y,z)$ 和抓取姿态 $(pitch,yaw,roll)$。
\end{enumerate}

色块角度用于调整末端 yaw。程序采用 \texttt{angle\_gain=-1} 和方块角度折返函数，将图像中的平面旋转角映射到夹爪绕竖直方向的角度修正。

转换函数默认桌面高度为 \SI{80}{mm}，基础抓取姿态为 $(pitch,yaw,roll)=(0^\circ,-90^\circ,180^\circ)$。对于色块这种近似正方形目标，图像角度存在 $90^\circ$ 周期等价，因此程序使用角度折返函数将旋转修正限制在 $[-45^\circ,45^\circ)$ 内，再叠加到末端 yaw 上。这样可以避免同一个正方形目标因为 \texttt{minAreaRect} 返回角度跳变而导致夹爪姿态突变。

坐标转换是视觉抓取中最容易累积误差的环节。检测器输出的是二维像素中心，而机械臂需要的是三维空间中的抓取点。程序通过桌面平面假设将三维问题简化为射线与平面的交点计算：只要物体位于同一工作平面附近，就可以用固定高度估计抓取位置。但如果物体高度明显变化，或相机与桌面之间的外参发生移动，转换结果就会偏离真实抓取点。

像素射线与桌面平面求交的实现如代码 \ref{lst:exp5-convert-code} 所示。它先将像素点去畸变为相机归一化射线，再通过相机到机械臂坐标变换求出机器人坐标系下的桌面交点。

\lstinputlisting[
  language=Python,
  caption={像素点转换到机械臂桌面坐标的核心代码},
  label={lst:exp5-convert-code}
]{code/exp5_pixel_to_robot.py}

\subsection{抓取流程设计}

完整抓取流程由 \path{daran/color_pick_place/main.py} 实现。程序先构造检测器，再组织“识别、选择、规划、仿真、执行确认”的流程。颜色识别模式下，检测器为 \texttt{ColorBlockDetector}；程序会过滤已经在目标区域附近的物块，选择候选目标，然后构造 pick 和 place 两个 \texttt{RobotPose}。本组仅分析和仿真该流程，没有进入硬件执行。

抓取动作序列为：
\[
\text{pick\_up}\rightarrow \text{pick\_down}\rightarrow \text{pick\_up}\rightarrow
\text{transit}\rightarrow \text{place\_up}\rightarrow \text{place\_down}\rightarrow
\text{place\_up}.
\]
其中每个位姿都包含末端位置、姿态和夹爪开合状态。程序调用 \texttt{JointActionMotion.\allowbreak solve()} 将位姿动作转换为关节动作，并支持仿真预览；代码中的实机执行入口本次未使用。

程序默认 standby 关节角为 $[0,90,0,0,0,0]^\circ$，与前面实验保持一致。放置点按颜色区分：红色和黄色目标默认放在 $(180,172,90)$ mm，绿色目标放在 $(125,172,90)$ mm，蓝色目标放在 $(50,172,90)$ mm。为了避免重复搬运已经处于目标区的物块，程序设置了 \SI{40}{mm} 的目标区忽略半径；如果检测到的物块已经位于某个目标附近，则不再选择它作为待抓取对象。

抓取动作中的上方点由目标点高度加上 approach height 得到，过渡点位于较高的 transit z。这样做延续了实验三的 safe-grasp-safe 思路：视觉只负责改变抓取点的位置，运动执行仍然遵守“先上方、再下降、再抬升”的安全模板。也就是说，实验五不是把视觉结果直接发送给机械臂，而是把视觉结果嵌入一个已有的安全动作结构中。

\subsection{实验步骤与异常处理}

程序流程首先读取工作区图像，随后根据目标颜色进行分割和轮廓筛选。若检测到多个候选目标，程序会根据目标类型、位置和是否已经位于放置区域附近进行过滤。选定候选目标后，系统将像素中心转换为机械臂坐标，生成抓取上方点、抓取点、过渡点、放置上方点和放置点，最后生成仿真预览。本组没有继续执行真实动作。

异常情况主要包括未检测到目标、检测到过小轮廓、目标位于不可达区域、坐标转换失败和逆运动学求解失败。对这些情况，合理处理方式不是强行执行，而是停止当前循环、提示操作者调整光照或物体位置，并重新识别。视觉抓取系统的安全性很大程度上来自这种“识别不可靠就不执行”的保守策略。

\subsection{离线图像与程序输出分析}

图 \ref{fig:task5-visual-grasp} 展示了已有图像上锁定候选目标的离线处理结果。该图只能说明程序能够在该样例中生成目标框和中心信息，不能证明实时相机检测稳定，也不能证明机械臂能够抓取该目标。

\begin{figure}[H]
  \centering
  \includegraphics[width=0.72\textwidth]{task5_visual_grasp_locked.jpg}
  \caption{颜色视觉抓取目标锁定结果}
  \label{fig:task5-visual-grasp}
\end{figure}

实验中还通过人工示教建立了像素到机械臂坐标的经验仿射关系，\path{calib_result.json} 中记录的二维仿射矩阵为：
\[
\begin{bmatrix}
x_r\\y_r
\end{bmatrix}
=
\begin{bmatrix}
0.0146&0.2952&140.9101\\
0.2440&0.0154&-224.5123
\end{bmatrix}
\begin{bmatrix}
u\\v\\1
\end{bmatrix},
\]
其文件中记录的 RMS 拟合误差约为 \SI{14.10}{mm}。该数值描述已有标定点上的二维拟合残差，不等同于真实机械臂末端到达误差，且本组未通过实物抓取检验该映射。

该仿射标定共使用 12 条记录，每条记录同时保存目标颜色、机器人坐标、相机像素中心和色块角度。与实验四中的针孔模型相比，二维仿射标定更依赖示教数据本身，优点是实现简单并能吸收一部分系统性偏差，缺点是难以解释相机姿态、镜头畸变和高度变化的物理来源。对于固定相机和固定桌面的课程实验，它可以作为理论标定之外的经验补偿。

从代码结构看，颜色识别方法实现简单、计算量较小，不需要训练模型；但阈值会受到光源、背景和相机曝光参数影响。由于本组没有进行实时场景和实物抓取测试，不能判断当前阈值是否足以支持真实闭环，只能将其作为后续测试的初始参数。

本实验代码体现了视觉规划相对于固定参考点回放的潜在优势：程序可以根据图像中的目标位置重新生成 pick pose 和末端 yaw。但这一能力是否能在物体平移、旋转和光照变化下可靠工作尚未验证。此外，现有视觉识别只发生在抓取前，执行过程中没有持续视觉反馈。

\subsection{调参与鲁棒性讨论}

颜色识别方法的调参主要集中在四个方面：HSV 阈值、面积阈值、长宽比范围和形态学核大小。HSV 阈值决定哪些像素会被认为属于目标颜色；面积阈值用于过滤噪声和小反光区域；长宽比范围用于保留近似方块形状的目标；形态学核大小则影响 mask 的连通性。参数过宽会引入误检，参数过窄会导致漏检，因此需要在当前光照、背景和相机曝光条件下折中。

蓝色目标在代码中进行了额外处理：程序先对图像进行简单白平衡和饱和度增强，再在 HSV 阈值外叠加 BGR 通道约束，要求蓝色通道明显高于红色和绿色通道。绿色目标也使用类似的通道约束。这些设置是针对可能出现的暗光和反光问题设计的规则，其稳定性仍需通过不同光照条件下的图像测试确认。

颜色视觉抓取的鲁棒性边界也比较清楚：当目标颜色与背景差异明显、光照稳定、目标不遮挡时，该方法速度快、实现简单；当光照变化大、背景中存在相近颜色、目标表面反光或多个目标相互遮挡时，检测结果会变得不稳定。若误检目标进入坐标转换和运动规划，机械臂可能抓向错误位置。因此系统必须在检测不到目标、坐标转换失败或 IK 求解失败时停止，而不是继续执行。

从工程结构看，颜色识别适合作为视觉抓取代码的第一版前端，因为它可以连接“图像、坐标、规划、执行接口”各模块。实验六在此结构上保留坐标转换和抓取后端，只替换目标发现方式；两种方案的实际效果都需后续验证。

\subsection{小结}

基于颜色识别的程序完成了从图像目标、坐标转换到机械臂动作规划接口的代码连接，并生成了仿真所需的候选位姿。由于没有进行实时相机和真实机械臂验证，当前不能评价抓取成功率和环境适应性。实验六进一步在相同后端接口上引入 YOLO 检测代码。
